\chapter{神经通讯}\label{ch:g8:nervous-communication}

\Cref{ch:g5:brain-in-command}中画出的回路——报告、
决定、下令——并测得其耗时五分之一秒。本年
我们要打开线缆。沿\omterm{def:g8:nervous-communication:neuron}{神经}究竟传递着什么？
\omterm{def:g8:nervous-communication:neuron}{神经}由什么构成？而在一根导线终止、另一根导线
开始之处的间隙里又会发生什么——那些间隙，
事实将表明，记忆与许多毒物都在其中
发挥作用？

\section{布线及其细胞}

\begin{definition}[神经元与神经信息]\label{def:g8:nervous-communication:neuron}
\omterm{def:g5:brain-in-command:system}{神经系统}的工作\omterm{def:g5:first-look-at-cells:cell}{细胞}是
\emph{神经元}\index{神经元}
（\cref{rem:g6:cell-unit-of-life:twoways}已列出它们）：
每个神经元都有一个\omterm{def:g5:first-look-at-cells:cell}{细胞}体和一根细长的纤维——有些纤维
从\omterm{def:g4:classifying-animals:vertebrate}{脊柱}一直延伸到脚趾，是全身最长的\omterm{def:g5:first-look-at-cells:cell}{细胞}。一
\emph{神经}\index{神经}是由数以千计的
这样的纤维捆成的线缆。沿每根纤维传递的是
\emph{神经信息}\index{神经信息}：短暂的电信号，
以高达 \qty{100}{m/s} 的速度飞驰
（\cref{def:g5:brain-in-command:system}中的图，
此时补上了它的载体）。
\end{definition}

\begin{proposition}[按频率编码]\label{prop:g8:nervous-communication:code}
信号本身都是一样的——变化的是
其\emph{速率}。轻触沿感觉纤维上传的是缓慢的
轻叩；用力的一触则是急促的鼓点。\omterm{def:g5:brain-in-command:system}{神经系统}的
语言只有一个“词”，以不同速率重复：
强度以频率的形式编码，
在每条线路上，无论传入还是传出。
\end{proposition}
\admitted

\begin{example}[三条线路，一种语言]\label{ex:g8:nervous-communication:lines}
回路中的各角色（\cref{prop:g5:brain-in-command:loop}），
经过改接：\emph{感觉\omterm{def:g8:nervous-communication:neuron}{神经元}}负责传递报告——
从\omterm{def:g1:the-five-senses:sense}{皮肤}到\omterm{def:g5:brain-in-command:system}{脊髓}，从眼睛到脑；\emph{运动\omterm{def:g8:nervous-communication:neuron}{神经元}}
负责传递命令——从\omterm{def:g5:brain-in-command:system}{脊髓}到\omterm{def:g3:bones-and-muscles:muscle}{肌肉}，止于纤维与\omterm{def:g3:bones-and-muscles:muscle}{肌肉}相接之处，
其急促的鼓点决定了
收缩的强度
（\cref{def:g3:bones-and-muscles:muscle}）；而在两者之间，
即脑与\omterm{def:g5:brain-in-command:system}{脊髓}之中，数量难以计数的
\emph{中间\omterm{def:g8:nervous-communication:neuron}{神经元}}负责做出决定。无论反射还是交响，
运行的都是同样的信号，
只是在不同线路上速率各异。
\end{example}

\section{间隙：突触}

\begin{definition}[突触]\label{def:g8:nervous-communication:synapse}
\omterm{def:g8:nervous-communication:neuron}{神经元}彼此并不接触。一根纤维的末端与
下一个\omterm{def:g5:first-look-at-cells:cell}{细胞}相接之处，存在一道微观的间隙——
\emph{突触}\index{突触}。电信号
无法跨越它；取而代之的是，抵达的纤维向间隙中
释放\emph{化学信使}，它穿过间隙，
在对岸与受体结合；如果到达的量足够，
便在下一个\omterm{def:g8:nervous-communication:neuron}{神经元}中重新产生信号。电信号
沿导线传递，化学信号则出现在每个连接处：整个
系统交替使用这两种方式。
\end{definition}

\begin{proposition}[为何是间隙而非焊点]\label{prop:g8:nervous-communication:whygaps}
间隙看似缺陷，实则是这一系统的精妙之处：
\begin{itemize}
  \item 使连接\emph{可调}——经常被穿越的\omterm{def:g8:nervous-communication:synapse}{突触}会变得更强，
        闲置的\omterm{def:g8:nervous-communication:synapse}{突触}则会衰退：
        这便是学习
        （\cref{rem:g5:brain-in-command:training}中的
        “更好的路径”，正是被加强的
        \omterm{def:g8:nervous-communication:synapse}{突触}）；
  \item 使信号\emph{单向}——信使
        只在一侧释放，因而信号无法
        逆向传递；
  \item 使一个\omterm{def:g8:nervous-communication:neuron}{神经元}能\emph{权衡}成千上万的
        输入——兴奋性与阻尼性的信号在每个\omterm{def:g5:first-look-at-cells:cell}{细胞}
        发放之前相加：这正是回路中的决定，
        在连接处完成。
\end{itemize}
\end{proposition}
\admitted

\begin{example}[杂耍者的突触]\label{ex:g8:nervous-communication:juggler}
\Cref{rem:g5:brain-in-command:training}中的杂耍者，
终于得到解释：一个月的练习使相同的
信号经由相同的连接传导，直到这些\omterm{def:g8:nervous-communication:synapse}{突触}
得到加强——原先要费力权衡的路由
如今已预先完成权衡。技能就是在间隙处固化下来的化学变化；
“熟能生巧”这句俗语，
说的其实是\omterm{def:g8:nervous-communication:synapse}{突触}的事。
\end{example}

\begin{omfigure}
\includegraphics[width=0.44\linewidth]{images/book1/ai/g8-nerve-juggler.jpg}

{\small 三只球自己飞舞：一个月里加强的
\omterm{def:g8:nervous-communication:synapse}{突触}，已被拍摄下来。}
\end{omfigure}

\begin{remark}[间隙是脆弱之处]\label{rem:g8:nervous-communication:vulnerable}
凡从外部作用于\omterm{def:g5:brain-in-command:system}{神经系统}的，皆作用于
\omterm{def:g8:nervous-communication:synapse}{突触}——这是信号变为化学物质的唯一位置，
可对接、可仿冒。麻醉剂使它们沉寂；
毒液则使其堵塞（有些通过阻断
\omterm{def:g3:bones-and-muscles:muscle}{肌肉}连接处的信使致人麻痹，有些则使其泛滥）；而
\cref{prop:g5:healthy-choices:substances}中的酒精，以及
一切改变情绪或感知的药物，都是通过
模拟、阻断或耗尽间隙中的化学信使来起作用。
\Cref{ch:g8:protecting-the-nervous-system}正是建立在
这一点之上。
\end{remark}

\section{快速电线，慢速邮件：两种系统比较}

\begin{proposition}[神经与激素的对比]\label{prop:g8:nervous-communication:compare}
人体现在展示出两种通信系统，它们
各自承担截然相反的任务：
\begin{itemize}
  \item \emph{\omterm{def:g8:nervous-communication:neuron}{神经}}：信号以毫秒计，沿专用的
        线路传至精确的目标——一块
        \omterm{def:g3:bones-and-muscles:muscle}{肌肉}或一个腺体——信号停止时
        作用即消：这是接住物体、说话与
        反射的系统；
  \item \emph{\omterm{def:g8:puberty:hormones}{激素}}
        （\cref{def:g8:puberty:hormones}）：信使以
        分钟至数月计，经\omterm{prop:g4:journey-of-food:absorb}{血液}广播至每一个
        能够接收信号的器官，指令缓慢、遍及全身的
        程序——生长、循环，以及\omterm{def:g8:puberty:puberty}{青春期}的重建。
\end{itemize}
电报与邮政：人体同时运行这两者，
\cref{ch:g8:hormonal-communication}专门讲
解邮政。
\end{proposition}
\admitted

\begin{method}[按系统诊断一种通信]\label{met:g8:nervous-communication:diagnose}
对任一身体协调活动，判定其所属系统：
\begin{enumerate}
  \item 计时：以毫秒计——\omterm{def:g8:nervous-communication:neuron}{神经}；以分钟至
        数年计——\omterm{def:g8:puberty:hormones}{激素}（或两者接力）；
  \item 测其作用范围：作用于一个精确地址——\omterm{def:g8:nervous-communication:neuron}{神经}；
        遍及全身——\omterm{def:g8:puberty:hormones}{激素}；
  \item 查明路径：可指明的导线通路——\omterm{def:g8:nervous-communication:neuron}{神经}；
        经\omterm{prop:g4:journey-of-food:absorb}{血液}广播——\omterm{def:g8:puberty:hormones}{激素}；
  \item 检查持续性：信号消失即止——
        \omterm{def:g8:nervous-communication:neuron}{神经}；信使仍在循环时持续——
        \omterm{def:g8:puberty:hormones}{激素}。
\end{enumerate}
\end{method}

\begin{example}[方法在三个实例中的应用]\label{ex:g8:nervous-communication:cases}
一阵气流吹来时的眨眼：毫秒级，作用于一圈
\omterm{def:g3:bones-and-muscles:muscle}{肌肉}，走导线通路，瞬间即逝——\omterm{def:g8:nervous-communication:neuron}{神经}。生长
骤增：以年计，遍及全身，经\omterm{prop:g4:journey-of-food:absorb}{血液}，
持续数季——\omterm{def:g8:puberty:hormones}{激素}。短跑中加速的心跳
（\cref{prop:g7:organs-at-work:supply}）：两者接力——
\omterm{def:g8:nervous-communication:neuron}{神经}信号在一次搏动内即提速，而
经\omterm{prop:g4:journey-of-food:absorb}{血液}的信使（即受惊时你感到的那阵著名的
热浪）使全身的警觉状态持续
数分钟。先是电报，邮政随后。
\end{example}

\section{练习}

\begin{exercise}[$\star$]\label{exo:g8:nervous-communication:1}
什么是\omterm{def:g8:nervous-communication:neuron}{神经元}？以\omterm{def:g8:nervous-communication:neuron}{神经元}而言，什么是\omterm{def:g8:nervous-communication:neuron}{神经}？
\end{exercise}

\begin{exercise}[$\star$]\label{exo:g8:nervous-communication:2}
什么沿纤维传导，速度多快——
强度又如何编码？
\end{exercise}

\begin{exercise}[$\star$]\label{exo:g8:nervous-communication:3}
说出回路中的三类\omterm{def:g8:nervous-communication:neuron}{神经元}，并指出各自的
通路。
\end{exercise}

\begin{exercise}[$\star$]\label{exo:g8:nervous-communication:4}
什么是\omterm{def:g8:nervous-communication:synapse}{突触}？什么物质穿过它？
\end{exercise}

\begin{exercise}[$\star$]\label{exo:g8:nervous-communication:5}
依据\cref{prop:g8:nervous-communication:whygaps}，
说出间隙提供的三项服务。
\end{exercise}

\begin{exercise}[$\star$]\label{exo:g8:nervous-communication:6}
比较\omterm{def:g8:nervous-communication:neuron}{神经}与\omterm{def:g8:puberty:hormones}{激素}在速度、作用范围、路径
及持续性上的差异。
\end{exercise}

\begin{exercise}[$\star\star$]\label{exo:g8:nervous-communication:7}
用本章的词汇重写
\cref{ex:g5:brain-in-command:catch}中的接球过程：
说明各类\omterm{def:g8:nervous-communication:neuron}{神经元}，并指出\omterm{def:g8:nervous-communication:synapse}{突触}在何处
进行信号的权衡。
\end{exercise}

\begin{exercise}[$\star\star$]\label{exo:g8:nervous-communication:8}
“练习使之永久。”请结合间隙解释这一口号，
并以杂耍者为例。
\end{exercise}

\begin{exercise}[$\star\star$]\label{exo:g8:nervous-communication:9}
为何一切从外部作用于\omterm{def:g5:brain-in-command:system}{神经系统}的事物
都作用于\omterm{def:g8:nervous-communication:synapse}{突触}？\omterm{def:g8:nervous-communication:synapse}{突触}有何特别之处？
\end{exercise}

\begin{exercise}[$\star\star$]\label{exo:g8:nervous-communication:10}
对下列情形运用\cref{met:g8:nervous-communication:diagnose}：（a）
从热锅边缩手；（b）\omterm{def:g8:puberty:puberty}{青春期}的嗓音变化；
（c）惊吓后仍持续的热浪。
\end{exercise}

\begin{exercise}[$\star\star$]\label{exo:g8:nervous-communication:11}
一种毒液阻断了\omterm{def:g3:bones-and-muscles:muscle}{肌肉}连接处的信使。
预测受害者的状态，并指出回路中被切断的环节。
\end{exercise}

\begin{exercise}[$\star\star\star$]\label{exo:g8:nervous-communication:12}
“沿导线为电信号，在连接处为化学信号——
一切有趣的事都发生在连接处。”
请从三个方面论证最后一句：学习、
决策与脆弱性。
\end{exercise}

\section{问题：反射实验台}

\begin{problem}[{周末问题——膝跳反射，
加装仪器}]\label{pb:g8:nervous-communication:1}
经典的台上演示：坐着的志愿者，
膝盖骨下方的小锤，足部的不自主
踢腿——现在用计时器加以测量，并结合本章内容。

\medskip\noindent\textbf{第一部分——回路。}
\begin{enumerate}
  \item 小锤拉伸肌腱，
        足部约在五十分之一秒后踢出。
        请列出回路：涉及哪些类型的\omterm{def:g8:nervous-communication:neuron}{神经元}，经由何处
        ——以及脑为何不在这一回路中。
  \item 五十分之一秒比
        接球的五分之一秒要快。
        请结合回路的长度与连接数，解释这一差异。
  \item 志愿者即使事先得到警告，仍无法抑制
        这一踢腿。这说明反射在何处被决定——
        以及\omterm{def:g8:nervous-communication:synapse}{突触}处的单向传递，
        说明了什么？
  \item 踢腿之后，志愿者
        \emph{感到}轻敲。请解释这一延迟的
        第二段较慢的传导路径。
\end{enumerate}

\medskip\noindent\textbf{第二部分——实验台变式。}
\begin{enumerate}[resume]
  \item 轻敲更用力，踢腿随之增强。
        在\cref{prop:g8:nervous-communication:code}中，
        哪条编码规则得到了体现？
  \item 校医指出，这一反射
        每次体检都会检查。能够引出的
        对称膝跳反射能说明\omterm{def:g5:brain-in-command:system}{脊髓}、\omterm{def:g8:nervous-communication:neuron}{神经}
        与连接处的何种状况？
  \item 咬紧下颌的志愿者踢得更重——
        这是一种已知的增强现象。这说明
        中间\omterm{def:g8:nervous-communication:neuron}{神经元}的信号竟也能抵达
        “无脑”回路？
  \item 腿部局部麻醉会同时消除踢腿
        与感觉。请依据
        \cref{rem:g8:nervous-communication:vulnerable}指出作用位点。
\end{enumerate}

\medskip\noindent\textbf{第三部分——书面报告。}
\begin{enumerate}[resume]
  \item 将该反射画成
        \cref{prop:g5:brain-in-command:loop}中的三步
        回路，并标出此反射在何处走了捷径。
  \item 用两列表格比较反射与接球：
        路径、连接处、时间、可抑制性、
        练习所起的作用。
  \item 报告的最后一节问道：“为何快速的保护性
        回路要绕过\omterm{def:g5:brain-in-command:system}{大脑}？”请
        运用热锅的逻辑作答
        （见\cref{exo:g5:brain-in-command:5}）。
  \item 用一句话收尾：本实验台同时展示了
        \omterm{def:g5:brain-in-command:system}{神经系统}的语言、线路与
        连接处。
\end{enumerate}
\end{problem}
